% ============================================================
%  ch5 Implementation and Optimization.tex
%  Chapter 5: Implementation and Optimization
%  Master's Thesis -- SQL MATCH_RECOGNIZE on Pandas DataFrames
% ============================================================
% Provide \mr if not already defined by the main document
\providecommand{\mr}{\texttt{MATCH\_RECOGNIZE}}
% !TEX root = Thesis.tex
\chapter{Implementation and Optimization}
\label{chap:implementation}

This chapter shows how the validated logical plan from
Chapter~\ref{chap:methodology} becomes an executable query.  A pattern
is compiled into a finite automaton and applied to ordered
\texttt{pandas} partitions.  Guarded fast paths reduce repeated Python
work without changing the supported semantics.
Section~\ref{sec:meth-automata} covers automata compilation,
Section~\ref{sec:meth-pandas} the DataFrame execution model, and
Section~\ref{sec:bg-performance} the optimization layer.

% ----------------------------------------------------------------
\section{Pattern Compilation with Finite Automata}
\label{sec:meth-automata}
% ----------------------------------------------------------------

Row-pattern matching is a sequence-recognition problem.  A row may
satisfy one or more pattern variables.  A match assigns consecutive
rows to those variables in an order accepted by the pattern.  Finite
automata provide a natural representation for this ordered
recognition task~\cite{HopcroftUllman}.  The engine therefore compiles
the pattern before it scans the data.

\paragraph{From pattern to NFA.}
The first automaton is a \emph{non-deterministic finite automaton}
(NFA).  It may follow several branches and take
\emph{$\varepsilon$-transitions}, which consume no row.  These
properties represent alternation, optional parts, and repetition
directly.  The builder follows Thompson's construction.  Each pattern
part becomes a small fragment, and the fragments are joined to form the
whole pattern~\cite{Thompson1968,DragonBook}.  The implementation also
carries SQL-specific metadata for conditions, exclusions, anchors, and
preference order.  Figure~\ref{fig:meth-nfa} shows a simplified NFA for
$A^{+}\,B?$ that accepts the same row sequences as the compiled
pattern.  Intermediate wrapper states used during construction are not
shown.

\begin{figure}[htbp]
  \centering
  \begin{tikzpicture}[
      x=1cm,
      y=1cm,
      nfastate/.style={circle, draw=black!80, line width=0.85pt,
                       fill=white, minimum size=0.98cm, inner sep=0pt,
                       font=\small},
      nfastart/.style={nfastate, fill=blue!7},
      nfaaccept/.style={nfastate, double, double distance=1.25pt,
                        fill=green!7},
      nfaconnector/.style={circle, draw=black!65, line width=0.7pt,
                           fill=black!3, minimum size=0.58cm,
                           inner sep=0pt, font=\scriptsize},
      entry/.style={-{Stealth[length=2.1mm]}, line width=0.75pt,
                     draw=black!75},
      consume/.style={-{Stealth[length=2.2mm]}, line width=1pt,
                       draw=blue!55!black},
      emptyedge/.style={-{Stealth[length=2.2mm]}, line width=0.8pt,
                         densely dashed, draw=black!65},
      edgelabel/.style={font=\small, fill=white, inner sep=1.3pt},
      fragmentline/.style={draw=black!45, line width=0.55pt},
      fragmenttext/.style={font=\footnotesize, text=black!85,
                            align=center},
      role/.style={font=\scriptsize, text=black!72, align=center},
      keytext/.style={font=\scriptsize, text=black!78, align=left}
    ]
    \node[nfastart]    (q0) at (0,0)    {$q_0$};
    \node[nfaaccept]   (q1) at (3.15,0) {$q_1$};
    \node[nfaconnector](r)  at (6.60,0) {$r$};
    \node[nfaaccept]   (q2) at (9.75,0) {$q_2$};

    % Thin bands identify the two pattern fragments without enclosing
    % the automaton in decorative panels.
    \draw[fragmentline] (-0.35,2.10) -- (3.50,2.10);
    \draw[fragmentline] (-0.35,2.00) -- (-0.35,2.20);
    \draw[fragmentline] (3.50,2.00) -- (3.50,2.20);
    \node[fragmenttext] at (1.58,2.43)
      {$A^{+}$ \enspace one or more \texttt{A} rows};
    \draw[fragmentline] (3.65,2.10) -- (10.10,2.10);
    \draw[fragmentline] (3.65,2.00) -- (3.65,2.20);
    \draw[fragmentline] (10.10,2.00) -- (10.10,2.20);
    \node[fragmenttext] at (6.88,2.43)
      {$B?$ \enspace zero or one \texttt{B} row};

    \draw[entry] (-1.25,0) -- (q0);
    \draw[consume] (q0) -- node[edgelabel, above]{\texttt{A}} (q1);
    \path[consume]
      (q1) edge[loop above, looseness=5, min distance=11mm]
      node[edgelabel]{\texttt{A}} (q1);
    \draw[emptyedge] (q1) -- node[edgelabel, above]{$\varepsilon$} (r);
    \draw[consume] (r) -- node[edgelabel, above]{\texttt{B}} (q2);

    \node[role] at (0,-1.30)    {initial};
    \node[role] at (3.15,-1.30) {accept: zero\\\texttt{B} rows};
    \node[role] at (6.60,-1.30) {ready for\\\texttt{B}};
    \node[role] at (9.75,-1.30) {accept: one\\\texttt{B} row};

    % A compact, unboxed key keeps the graph readable in grayscale.
    \draw[fragmentline] (-1.25,-1.93) -- (10.75,-1.93);
    \node[keytext, font=\scriptsize\bfseries, anchor=west]
      at (-1.20,-2.28) {Key};
    \draw[consume] (-0.35,-2.28) -- (0.43,-2.28);
    \node[keytext, anchor=west] at (0.55,-2.28)
      {consumes one matching row};
    \draw[emptyedge] (4.25,-2.28) -- (5.03,-2.28);
    \node[keytext, anchor=west] at (5.15,-2.28)
      {$\varepsilon$: consumes no row};
    \node[nfaaccept, minimum size=0.44cm] at (9.25,-2.28) {};
    \node[keytext, anchor=west] at (9.60,-2.28) {accept};
  \end{tikzpicture}
  \caption[Language-equivalent NFA for $A^{+}\,B?$]{Language-equivalent NFA for $A^{+}\,B?$: $q_1$ accepts with zero
    \texttt{B} rows, $q_2$ after one.}
  \label{fig:meth-nfa}
\end{figure}

An \texttt{A} or \texttt{B} transition consumes one row for which that
variable's \texttt{DEFINE} predicate is true.  State $q_1$ is reached
after one or more \texttt{A} rows and is accepting because \texttt{B}
is optional.  From $q_1$, the NFA may accept, repeat \texttt{A}, or
enter the optional fragment through the $\varepsilon$-transition.  The
transition into the connector state $r$ consumes no row.  Consuming one
\texttt{B} row then reaches the second accepting state, $q_2$.  These
alternatives show nondeterminism.  The default greedy preference is
stored as transition metadata and is not represented by the simplified
state graph.

\paragraph{Reading the pattern.}
A left-to-right tokenizer reads variables, alternation, grouping,
anchors, exclusions, and \texttt{PERMUTE}.  It normalizes
\texttt{*}, \texttt{+}, \texttt{?}, and
\texttt{\{$n$,$m$\}} into minimum and maximum counts together with a
greedy or reluctant preference.  The builder can then handle every
quantifier through one internal representation.

\paragraph{Building the automaton.}
Each token contributes an NFA fragment.  When several continuations are
possible, transition metadata preserves their pattern order.
\texttt{PERMUTE} over $n$ variables represents all $n!$ orders.
Construction is therefore guarded at eight variables and an estimated
50{,}000 branches.  Exceeding either guard raises an explicit
complexity error.  Rows inside \texttt{\{-\ \ldots\ -\}} still take
part in matching; only their output is excluded.

\paragraph{From NFA to DFA.}
The NFA is converted into a deterministic finite automaton (DFA) by
subset construction.  Each DFA state represents a complete set of NFA
states.  Its $\varepsilon$-closure includes every state reachable
without consuming a row.  Reusing a previously seen set prevents
duplicate DFA states~\cite{HopcroftUllman}.

The resulting automaton has no $\varepsilon$-transitions.  For eligible
row-local pattern structures, the matcher can advance through the
ordered partition without keeping the complete NFA frontier at every
row.  Predicates or structural choices that depend on tentative
assignments use the exact search described in
Section~\ref{subsec:meth-exec-backtracking}.

Subset construction is exponential in the worst case.  The engine
therefore applies a memory-aware state budget, an iteration limit, and
a structural ceiling of 50{,}000 states.  The memory budget may lower
that ceiling on a constrained system.  Reaching any limit raises a
construction error.  The engine never executes or caches a partial DFA,
and it never truncates an NFA subset.  None of the evaluated patterns
reached these limits.

After construction, unreachable states are removed and transitions
are normalized.  The current reduction is conservative; it does not
perform full DFA minimization.  Transition metadata is preserved so
the runtime can apply the required greedy and leftmost preference.

% ----------------------------------------------------------------
\section{Pandas DataFrame Execution Model}
\label{sec:meth-pandas}
% ----------------------------------------------------------------

\texttt{pandas}~\cite{McKinney2010} stores each DataFrame column as an
array.  This suits both row-by-row matching and the vectorized fast
paths in Section~\ref{sec:bg-performance}.  Figure~\ref{fig:system-architecture}
shows the system boundary, while Figure~\ref{fig:exec-pipeline} expands
Stage~5.  Its two inputs are the compiled execution plan and the
\texttt{pandas} DataFrame.  The SQL query is shown only as the upstream
source of that plan.  The figure presents five logical execution
steps: partitioning, ordering, match
enumeration with predicate evaluation, measure evaluation, and result
materialization.  An
optimized physical path may fuse some operations without changing this
logical flow.

\begin{figure}[htbp]
  \centering
  \begin{tikzpicture}[
      x=1cm,
      y=1cm,
      context/.style={draw=black!55, rounded corners=3pt,
                       line width=0.75pt, fill=black!3,
                       minimum height=0.82cm, align=center,
                       inner xsep=5pt, inner ysep=3pt,
                       font=\footnotesize},
      plan/.style={context, draw=blue!50!black, fill=blue!6,
                   line width=0.85pt},
      dataframe/.style={context, minimum width=1.95cm,
                         minimum height=1.55cm, inner sep=3pt},
      resultframe/.style={dataframe, draw=green!48!black,
                          fill=green!7, line width=0.85pt},
      station/.style={circle, draw=blue!58!black,
                       fill=blue!58!black, text=white,
                       line width=0.8pt, minimum size=0.66cm,
                       inner sep=0pt, font=\footnotesize\bfseries},
      operation/.style={font=\footnotesize, text=black!84,
                         text width=1.82cm, align=center, anchor=north},
      dataflow/.style={-{Stealth[length=2.2mm,width=1.55mm]},
                        draw=blue!55!black, line width=0.9pt},
      controlflow/.style={-{Stealth[length=1.8mm,width=1.3mm]},
                           draw=black!58, densely dashed,
                           line width=0.7pt},
      phasebar/.style={draw=black!42, line width=0.5pt},
      phasetext/.style={font=\footnotesize\bfseries,
                        text=black!70, fill=white, inner xsep=7pt},
      stageheader/.style={font=\footnotesize\bfseries,
                          text=black!72, fill=white, inner xsep=6pt},
      upstreamlabel/.style={font=\footnotesize, text=black!65,
                            fill=white, inner sep=1pt}
    ]
    % Compile-time control lane.
    \node[context, minimum width=2.05cm] (sql) at (-0.10,3.02)
      {\textbf{SQL query}\\[-1pt]
       {\scriptsize\texttt{MATCH\_RECOGNIZE}}};
    \node[plan, text width=7.20cm] (compiled) at (6.55,3.02)
      {\textbf{Compiled execution plan}\\[-1pt]
       {\footnotesize partition/order keys $\cdot$ \texttt{DEFINE}
        $\cdot$ compiled pattern\\[-1pt]
        measures $\cdot$ skip policy $\cdot$ output mode}};
    \draw[dataflow] (sql) --
      node[upstreamlabel, above] {Stages 1--4} (compiled);

    % Runtime data lane.
    \draw[phasebar] (1.25,1.56) -- (10.75,1.56);
    \node[stageheader] (exec) at (6.00,1.56)
      {Stage 5 --- DataFrame execution};
    \draw[controlflow] (compiled.south) -- (exec.north);

    \node[dataframe] (data) at (0,0)
      {\textbf{Input}\\[-1pt]
       {\scriptsize\texttt{pandas} DataFrame}\\[3pt]
       {\scriptsize
       \begin{tabular}{@{}ccc@{}}
         \texttt{key} & \texttt{t} & \texttt{value} \\
         \hline
         $\cdots$ & $\cdots$ & $\cdots$
       \end{tabular}}};

    \node[station] (s1) at (2.10,0) {1};
    \node[station] (s2) at (4.00,0) {2};
    \node[station] (s3) at (5.95,0) {3};
    \node[station] (s4) at (7.95,0) {4};
    \node[station] (s5) at (9.85,0) {5};

    \node[operation] at (2.10,-0.49)
      {\textbf{Partition rows}\\[-1pt]
       {\scriptsize\texttt{PARTITION BY}}};
    \node[operation] at (4.00,-0.49)
      {\textbf{Order rows}\\[-1pt]
       {\scriptsize\texttt{ORDER BY}}};
    \node[operation, text width=2.05cm] at (5.95,-0.49)
      {\textbf{Enumerate matches}\\[-1pt]
       {\scriptsize\texttt{DEFINE}\\[-1pt] preference and skip}};
    \node[operation, text width=1.90cm] at (7.95,-0.49)
      {\textbf{Evaluate measures}\\[-1pt]
       {\scriptsize\texttt{RUNNING} or \texttt{FINAL}}};
    \node[operation, text width=1.90cm] at (9.85,-0.49)
      {\textbf{Materialize result}\\[-1pt]
       {\scriptsize row mode\\[-1pt] and projection}};

    % Thin phase brackets group the five operations without extra panels.
    \draw[phasebar] (1.45,0.88) -- (4.65,0.88);
    \node[phasetext] at (3.05,0.88) {Prepare rows};
    \draw[phasebar] (5.05,0.88) -- (6.85,0.88);
    \node[phasetext] at (5.95,0.88) {Recognize matches};
    \draw[phasebar] (7.15,0.88) -- (10.65,0.88);
    \node[phasetext] at (8.90,0.88) {Build result};

    \node[resultframe] (result) at (12.00,0)
      {\textbf{Result}\\[-1pt]\textbf{DataFrame}\\[2pt]
       {\scriptsize
       \begin{tabular}{@{}cc@{}}
         \texttt{key} & \texttt{measure} \\
         \hline
         $\cdots$ & $\cdots$
       \end{tabular}}};

    % Solid arrows carry the DataFrame row flow.
    \draw[dataflow] (data) -- (s1);
    \draw[dataflow] (s1) -- (s2);
    \draw[dataflow] (s2) -- (s3);
    \draw[dataflow] (s3) -- (s4);
    \draw[dataflow] (s4) -- (s5);
    \draw[dataflow] (s5) -- (result);
  \end{tikzpicture}
  \caption{Inside the DataFrame execution stage.}
  \label{fig:exec-pipeline}
\end{figure}

\paragraph{Preparing the data.}
A match never crosses a partition boundary, so the DataFrame is first
split on the \texttt{PARTITION BY} columns---or treated as one partition
when there are none.  The engine checks the \texttt{ORDER BY} columns.
It keeps a partition that already follows the complete ordering
specification and sorts it otherwise.  Each key may use \texttt{ASC} or
\texttt{DESC} and its own null placement.  The sort is stable, so rows
with equal keys keep their DataFrame order.  This avoids repeating a
sort whose result is already present.  Because matches do not cross partitions,
the remaining steps can process one ordered partition at a time.

Each \texttt{DEFINE} condition is turned into a small function once, at
the start of the query.  A variable used in \texttt{PATTERN} but not in
\texttt{DEFINE} receives an always-true condition, as required by the
supported semantics~\cite{trino2022}.  Row-local conditions may be
evaluated for the whole partition.  Other conditions are evaluated only
when a candidate assignment needs them.

\paragraph{Finding the matches.}
Matching walks the rows in order.  At each position, it tests the
possible variable transitions and applies the pattern's preference.
Greedy quantifiers prefer a longer valid match.  Reluctant quantifiers
prefer a shorter one.  Alternation follows declaration order.  The
automata carry this ordering information, while the runtime checks that
the preferred path can still complete.

Each completed match records its start, end, variable assignments, and
excluded rows.  This state drives skip handling and measure evaluation.
Exclusion changes only the rows emitted; it does not change the rows
used by conditions or measures.

% ................................
\subsection{Exact Search for State-Dependent Conditions}
\label{subsec:meth-exec-backtracking}
% ................................

The deterministic walk is used only when the predicates and pattern
structure can be decided without revisiting tentative assignments.
Other conditions depend on the tentative match.  For example,
accepting a row changes a running average used by a later condition.
Some nested quantified alternatives also need a choice to be tested
before their preference can be decided.

These cases use a preference-ordered depth-first search with
backtracking.  The search tentatively assigns rows, evaluates the
affected conditions, and tries the next legal choice when a later step
fails.  It is iterative, so Python's recursion depth does not limit the
match length.  A state budget and a search budget bound the explored states and
steps.  If either is exhausted, the query raises an
explicit error instead of returning a partial or different match.
Eligible row-local patterns remain on deterministic fast paths.
Section~\ref{subsec:eval-worst-cases} measures one supported
state-dependent case.

% ................................
\subsection{Skip Modes, Output Modes, and Measures}
\label{subsec:meth-exec-output}
% ................................

After recording a match the engine must decide where to resume, which
controls whether matches can overlap.  SQL:2016 defines four
\texttt{AFTER MATCH SKIP} modes, summarized with their next-start
computation in Table~\ref{tab:skip-modes}.  Empty matches always advance
by one row.  A variable-based target that does not occur in the match,
or that points back to the current start, raises an error.  This
prevents a silent loop or a changed match sequence.

\begin{table}[htbp]
  \centering
  \caption{The four \texttt{AFTER MATCH SKIP} modes and the
    resulting next scan position}
  \label{tab:skip-modes}
  \begin{tabular}{lp{9cm}}
    \toprule
    \textbf{Mode} & \textbf{Next-start computation} \\
    \midrule
    \textsc{Past Last Row} (default) &
      $i \leftarrow \mathit{best.end} + 1$. No overlap is possible;
      the scan resumes immediately after the last matched row. \\[4pt]
    \textsc{To Next Row} &
      $i \leftarrow \mathit{best.start} + 1$. Overlapping matches
      are produced because only a single row is skipped regardless
      of match length. \\[4pt]
    \textsc{To First} $V$ &
      $i \leftarrow \min(\mathit{bindings}[V])$. The scan
      resumes at the \emph{first} row assigned to variable $V$
      ($V$ may also be a \texttt{SUBSET} alias). \\[4pt]
    \textsc{To Last} $V$ &
      $i \leftarrow \max(\mathit{bindings}[V])$. The scan
      resumes at the \emph{last} row assigned to variable $V$ (or
      \texttt{SUBSET} alias). \\
    \bottomrule
  \end{tabular}
\end{table}

The engine supports \texttt{ONE ROW PER MATCH} and
\texttt{ALL ROWS PER MATCH}.  The latter includes
\textsc{Show Empty Matches}, \textsc{Omit Empty Matches}, and
\textsc{With Unmatched Rows}~\cite{trino2022}.  Empty matches are
shown by default; omitting them removes their output rows, while
unmatched mode also emits input rows not covered by a non-empty match.
\texttt{CLASSIFIER()} can expose the variable assigned to an output
row.  Table~\ref{tab:rows-per-match} summarizes the row counts and
measure contexts.

\begin{table}[htbp]
  \centering
  \caption{Output modes and their row-count and measure-semantics
    behavior}
  \label{tab:rows-per-match}
  \begin{tabularx}{\textwidth}{>{\raggedright\arraybackslash}p{3.0cm}%
    >{\raggedright\arraybackslash}p{3.5cm}X}
    \toprule
    \textbf{Mode} & \textbf{Rows out} & \textbf{Measure evaluation context} \\
    \midrule
    \textsc{One Row Per Match}
      & 1 per match
      & Measures use the completed variable bindings. \\[6pt]
    \textsc{All Rows Per Match}
      & One per visible matched row; \textsc{Show Empty} adds one row
        for each empty match
      & Measures use their resolved \textsc{Running} or
        \textsc{Final} semantics for each emitted row. \\[6pt]
    \textsc{All Rows \ldots With Unmatched}
      & Visible matched rows, plus each unmatched input row; empty
        matches follow the selected empty-match option
      & Matched rows carry measure values.  Unmatched rows emit
        \texttt{NULL} for the measure columns. \\
    \bottomrule
  \end{tabularx}
\end{table}

Measure semantics are resolved before matching.  Common navigation and
aggregate forms are compiled once; other supported expressions use the
general evaluator.  A qualified aggregate uses rows assigned to its
variable or \texttt{SUBSET}.  An unqualified aggregate uses the
applicable \textsc{Running} match prefix or the complete \textsc{Final}
match.
Direct navigation plans use bounds-checked lookups into the ordered
partition, and an out-of-range reference returns \texttt{NULL}.
\texttt{MATCH\_NUMBER()} numbers matches within each partition, while
\texttt{CLASSIFIER()} reports the variable assigned to the output row.
Under \texttt{ALL ROWS PER MATCH}, a \textsc{Running} measure sees the
match prefix through the current row.  A \textsc{Final} measure sees
the complete match.

% ----------------------------------------------------------------
\section{Performance Optimization Techniques}
\label{sec:bg-performance}
% ----------------------------------------------------------------

One rule governs this whole layer: a faster path may never change the
answer.  Each technique below is therefore switched on only when the
engine can be sure it applies, and the ordinary path takes over
otherwise.  Differential tests compare eligible fast paths with the
general path on the same input and output mode.

The techniques remove repeated predicate work, avoid unnecessary
sorting and materialization, and reuse compiled structures.
Table~\ref{tab:perf-summary} gives their activation boundaries.

\paragraph{Deciding conditions ahead of time.}
The main cost in the general path is repeated condition evaluation.
Predicates that depend only on the current row, literals, and eligible
fixed-offset \texttt{PREV}/\texttt{NEXT} forms can be evaluated as
column operations.  Their Boolean results are stored once per
partition.  The matcher then reads a mask instead of calling the
expression evaluator for every candidate row.

The planner is conservative.  \texttt{FIRST}, \texttt{LAST},
\texttt{CLASSIFIER()}, aggregates, and cross-variable references depend
on the tentative match and cannot use a row-local result.  A mixed
\texttt{AND} predicate may still use its proven row-local conjuncts as
a rejection mask; the remaining conjuncts are evaluated exactly.
Supported expressions outside the compiled forms use the general
evaluator.  Unsupported forms are rejected.

The running-average condition used in the worst-case experiment
receives a separate
incremental plan.  The aggregate is updated as the matched prefix grows
instead of being recomputed from the beginning.  This removes the
quadratic repeated-aggregation cost for that supported form.  It does
not make every state-dependent predicate vectorizable
(Section~\ref{subsec:eval-worst-cases}).

\paragraph{Turning the search into a lookup.}
When all variable predicates have masks, the engine builds a compact
transition index for the DFA states.  Matching becomes array lookup and
integer comparison rather than repeated dictionary and function work.
Assignments can be stored as row segments instead of one Python object
per row.

A linear pattern such as \texttt{A B\{2,4\} C+} allows a second plan.
Run-length arrays record how many consecutive rows satisfy each
variable.  Failed token--position pairs are memoized, so the same
impossible suffix is not explored again.  Patterns with alternation use
the row-local DFA path when eligible.  State-dependent choices use the
exact matcher.  Neither case enters the linear run-length plan.

For eligible output modes, matching and materialization are fused.
One path emits \texttt{ONE ROW PER MATCH} results directly.  Another
publishes compact segments for \texttt{ALL ROWS PER MATCH}.  Both
require the default skip mode and fully compiled measures.  Any anchor,
empty-match rule, measure, or output option outside their proof
boundary returns to the complete path.

\paragraph{Computing measures once.}
Common measures are compiled once into direct plans.  Simple
\texttt{SUM}, \texttt{AVG}, \texttt{COUNT}, statistical aggregates,
and direct navigation can then use arrays or prefix state.  A
\texttt{RUNNING} aggregate over a long match is updated incrementally
instead of being recomputed for every output row.  The compiled layer
also covers selected arithmetic combinations and distinct aggregates.
Filters and expressions outside these proven forms use the general
measure evaluator.

\paragraph{Reducing repeated work and memory.}
Compiled automata are stored in a bounded, thread-safe
\emph{least-recently-used} (LRU) cache.  The key includes the pattern,
its \texttt{DEFINE} conditions, and compilation-relevant
\texttt{SUBSET} data.  Repeating the same structure therefore avoids
NFA and DFA construction.  Condition caches are also bounded.  For a
large partition, a small row sample warms common evaluation paths.

The executor avoids copying all grouped partitions before matching.
For eligible default-skip, non-\texttt{PERMUTE} queries, it exposes one
partition through a lazy row view with a bounded row cache.  Other
paths materialize partition rows and may retain them until output
assembly.  Match membership uses compact masks and segments where the
semantics allow it.  A conservative
final-output projection gathers only source columns needed by the
result; a wildcard or an expression it cannot prove safe keeps all
columns.  Eligible fast paths assemble
output columns directly before creating the result DataFrame.  These
choices lower temporary memory, but they do not make execution
streaming: it retains the input and result, and some paths also retain
materialized rows or whole-column working arrays.  A cache budget bounds cache capacity, a state
budget bounds automaton construction, and a search budget bounds exact
search.

\begin{table}[htbp]
  \centering
  \caption{Performance optimizations and their activation
    conditions}
  \label{tab:perf-summary}
  \begin{tabular}{p{3.8cm}p{4.2cm}p{4.4cm}}
    \toprule
    \textbf{Technique} & \textbf{Activation condition}
      & \textbf{Primary benefit} \\
    \midrule
    Ordered-input check
      & Requested \texttt{ORDER BY} sequence already present
      & Avoids an unnecessary sort and DataFrame copy \\[4pt]
    Vectorized condition masks
      & Row-local expression, or safe row-local conjuncts in
        a mixed \texttt{AND}
      & Replaces repeated Python evaluation with column operations
        and exact prefilters \\[4pt]
    Row-local DFA decision tables
      & All predicates are vectorized, with no empty-match, anchor,
        exclusion, \texttt{PERMUTE}, \texttt{SUBSET}, reluctant, or
        state-dependent requirement
      & Replaces repeated transition dispatch with compact indexing
        and segment assignments \\[4pt]
    Linear run-length plan
      & Default skip, a plain linear variable sequence, supported
        quantifiers, and row-local masks
      & Skips complete runs and memoizes failed suffixes \\[4pt]
    Fused output drivers
      & Default skip, no unmatched rows, anchors, or empty matches, an
        eligible matcher, and fully compiled measures
      & Avoids intermediate match and row dictionaries \\[4pt]
    Compiled measure and prefix plans
      & Supported direct navigation or aggregate form
      & Reuses prepared array state; avoids repeated prefix
        aggregation \\[4pt]
    LRU pattern cache
      & Caching enabled and the complete automaton fits the cache budget
      & Eliminates redundant NFA/DFA construction across repeated
        queries \\[4pt]
    Lazy rows and output projection
      & Partition path supports lazy access; selected source columns
        can be resolved safely
      & Reduces row copies, retained Python objects, and output width \\
    \bottomrule
  \end{tabular}
\end{table}

% ----------------------------------------------------------------
\section{Chapter Summary}
\label{sec:impl-summary}
% ----------------------------------------------------------------

This chapter followed a query from a compiled pattern to a result.  The
pattern becomes an NFA and then a complete, guarded DFA.  The input is
partitioned and ordered before matching.  Row-local queries use the
deterministic path only when their predicates and structure are
eligible.  Predicates or structural choices that depend on tentative
assignments use exact search and fail explicitly if its search budget
is exhausted.

Major speed gains come from safe column-wise condition evaluation,
linear run-length plans, compiled measures, fused output, and pattern
caching.  Lazy row access,
projection, compact masks, and the cache and state budgets reduce
temporary memory.  Every fast path is guarded.  If its conditions are
not met, the complete path remains responsible for the result.
Chapter~\ref{chap:evaluation} measures the resulting correctness,
performance, scalability, and memory behavior.
